\documentclass[11pt,a4paper]{article}
% Shared preamble for RuPhO English problem statements (match T1 2026 style).
% Use: \input{latex/rupho_en_preamble.tex} after \documentclass.
\usepackage[margin=1in]{geometry}
\usepackage{amsmath,amssymb}
\usepackage{graphicx}
\usepackage{float}
\usepackage{enumitem}
\usepackage{microtype}
\usepackage{caption}
\usepackage{tabularx}
\usepackage{hyperref}

\captionsetup{font=small,labelfont=bf,skip=6pt}

\setlist[itemize]{leftmargin=1.2cm}
\setlist[enumerate]{leftmargin=1.2cm}

% One outer box per subproblem: [id | statement | points] (no inner rules)
\newcommand{\problembox}[3]{%
  \par\addvspace{0.42em}%
  \noindent
  \setlength{\fboxsep}{7.5pt}%
  \setlength{\fboxrule}{0.95pt}%
  \fbox{%
    \begin{minipage}{\dimexpr\linewidth-2\fboxsep-2\fboxrule\relax}
    \begin{tabularx}{\linewidth}{@{}>{\raggedright\arraybackslash\bfseries}p{2.1em} @{\hspace{0.9em}} >{\raggedright\arraybackslash}X @{\hspace{0.9em}} >{\raggedleft\arraybackslash\bfseries}p{2.4em}@{}}
    #1 & #2 & #3
    \end{tabularx}
    \end{minipage}%
  }%
  \par\addvspace{0.32em}%
}

\title{\textbf{Crooked cylinder}}
\author{\normalsize X19 (2019) --- English translation}
\date{}
\begin{document}
\maketitle

A cylindrical vessel of height $H=50~\text{cm}$ with mirror walls is filled with a transparent liquid, the refractive index of which decreases with height according to a linear law. Near the bottom of the cylinder there is a point light source with power $P=100~\text{W}$. The refractive index of the liquid near the bottom is $n_1=1.50$, and at the surface $n_2=1.44$. The bottom of the vessel and its lid are covered with a substance that absorbs all radiation that hits it

\begin{figure}[H]
    \centering
    \includegraphics[width=0.72\textwidth]{T8_EN_assets/fig1.jpg}
    \caption{Figure 1.}
\end{figure}

\problembox{A1}{What radiation power $P_\text{}$ is absorbed by the bottom of the cylinder?}{0}

\problembox{A2}{Suppose a certain beam passes near the cylinder cover. What is the radius of curvature $R_\text{}$ of this ray?}{0}

\problembox{A3}{What is the radius of curvature $R_\text{}$ of the same ray directly near the source?}{0}

\vspace{1em}
\noindent\textit{Source:} \url{https://pho.rs/p/3059}

\end{document}